
\enteteExam{NFE106 -- Administration de bases de données}%
{Examen 2 septembre 2021 -- session 2 -- 3h}%
{\empty}%

\section{Stockage et indexation (6 pts)}

Soit l'extrait suivant de la base de données qui sert à la gestion d'une grande
 bibliothèque\,:

\begin{itemize}
  \item Livre (id, titre, id\_auteur, date\_publication, catégorie, éditeur)
  \item Emprunt (id\_personne, id\_livre, date\_début, date\_fin)
  \item Personne  (id, nom, prénom, âge)
\end{itemize}

\begin{enumerate}

\item Pour la table Livre, identifiez les clés primaires et étrangères.
\begin{solution}
Clé primaire: id\\
Clés étrangères: id\_auteur $\rightarrow$ AUTEUR(id)
\end{solution}

\item Donnez la commande SQL pour créer la table Livre. 
\begin{solution}
\begin{verbatim}
CREATE TABLE LIVRE(
    id NUMBER(10),
    titre VARCHAR2(50),
    id_auteur NUMBER(10),
    date_publication DATE,
    catégorie VARCHAR2(20),
    éditeur VARCHAR2(20),
    CONSTRAINT pk_livre PRIMARY_KEY id,
    CONSTRAINT fk_livre FOREIGN_KEY id_auteur REFERENCES AUTEUR(id)
\end{verbatim}
\end{solution}

\item On 
suppose que chaque enregistrement a une taille de 
250 octets et que chaque bloc a un entête de 150 octets. 
 Indiquez le nombre de blocs nécessaires si la table 
       contient 3 millions de livres. La taille d'un bloc est de 4\,000 octets.

\begin{solution}
Il y a 3 850 octets utiles dans un bloc. On place $3850/250=15$ enregistrements par bloc.  
Il faut 
        $\lceil \frac{3\,000\,000}{15} \rceil = 200\,000$ blocs.
\end{solution}

 \item On décide de trier la table Livre sur la date de publication et de construire
     un index non-dense à un niveau. Combien contient-t-il d'entrées, 
     Sachant qu'une adresse et une date occupent chacun 8 octets, combien d'entrées met-on dans
     un bloc et combien de blocs sont nécessaires\,?

\begin{solution}
Il y a 200 000 entrées de 16 octets chacune. . On place $4000/16=250$ entrées par bloc.  
Il faut 
        $\lceil \frac{200\,000}{250} \rceil = 800$ blocs.
\end{solution}

\end{enumerate}

\section{Indexation (4 points)}


On insère successivement les enregistrements suivants dans la table 
\texttt{Personne}, selon cet ordre :

\begin{center}
\begin{small}
\begin{tabular}{c|c|c|c}
\textbf{id} & \textbf{nom} & \textbf{prénom} & \textbf{âge} \\
\hline
\hline
1 & Martin & Georges & 30\\
\hline
2 & Dubois & Inna & 20 \\
\hline
3 & Thomas & Maxime & 14\\
\hline
4 & Richard & Boris & 45\\
\hline
5 & Petit & Gérald & 32\\
\hline
6 & Durand & Valérie & 43\\
\hline
7 & Leroy & Jacques & 34 \\
\hline
8 & Simon & Michel & 22 \\
\hline
9 & Laurent & Jean-Jacques & 63\\
\hline
10 & Moreau & Eddie & 12\\
\hline
11 & Morel & Philippe & 15\\
\hline
12 & Perrin & Gloria & 18 \\
\hline
13 & Clement & Bob & 19\\
\hline
14 & Dumont  &  Claire & 67\\
\hline
\end{tabular}
\end{small}
\end{center}

On met en place un index sur l'attribut \textbf{âge} de cette table. 

\begin{enumerate}

\item Construire l'arbre B d'ordre \textbf{2} (4 entrées par bloc)
correspondant à cet ensemble d'articles, en respectant l'ordre d'arrivée. 
Donner les étapes de construction intermédiaires importantes (celles qui produisent 
un changement de la structure de l'arbre).


\item On considère que cet arbre B  est stocké à raison d'un n{\oe}ud par bloc sur disque. On recherche les noms des personnes dont l'âge est entre 20 et 40 ans (inclus). 
Combien de blocs disque doivent être chargées au minimum pour répondre à cette requête en utilisant l'arbre B ? Justifier. 


\begin{solution}
Il s'agit ici d'une requête par intervalle: on commence par chercher les enregistrements pour lesquels l'âge est égal à 20 (lecture de 5 enregistrements). Il suffit ensuite d'exploiter le chaînage des feuilles pour trouver les autres enregistrements, en parcourant ici encore 5 enregistrements dans l'arbre B). On charge donc 10 
blocs d'index. Il faut ensuite récupérer les données pointées par l'index sur disque. Dans le meilleur cas (peu probable mais possible), toutes ces données sont stockées dans un même 
bloc. Donc on charge au minimum 11 blocs. 
\end{solution}

\end{enumerate}


\section{Optimisation (6 points)}

\begin{enumerate}

   \item Soit la requête suivante:
\begin{minted}{sql}
     select p.nom, p.prénom
     from Livre as l,Personne as p
     where l.titre = "Bilbo"
     and l.id_auteur = p.id
\end{minted}

  Donnez le plan d'exécution classique des boucles imbriquées indexées, en supposant
     que seules les clés primaires sont indexées.
     
    \item Que devient le plan si on ajoute un index sur le titre des livres 
     
\begin{minted}{sql}
     select e.date_debut, e.date_fin 
     from Livre as l, Emprunt as e
     where  l.titre = "Bilbo"
     and e.id_livre = l.id
\end{minted}

  
  \item En supposant qu'il n'y a pas d'index sur la table Livre, quelle serait la condition
   pour utiliser un index
   dans l'algorithme de jointure\,? Quel
   serait alors le plan d'exécution\,?  Quel plan serait possible sans index\,?

\end{enumerate}

\section{Concurrence (4 points)}

Soit l'exécution concurrente suivante

$$H = r_1[x] r_2[x] r_3[y] w_1[x] w_1[z] c_1 w_2[y] c_2 w_3[z] c_3$$

\begin{enumerate}

\item Donner la liste des conflits de H
\begin{solution}
    \begin{itemize}
        \item[sur x : ] $r_2[x] w_1[x]$
        \item[sur y : ] $r_3[y] w_2[y]$
        \item[sur z : ] $w_1[z] w_3[z]$
    \end{itemize}
\end{solution}

\item Donner le graphe de s\'erialisation de H. Que pouvez-vous déduire de ce
graphe? 
\begin{solution}
    $T_2 \rightarrow T_1, T_3 \rightarrow T_2, T_1 \rightarrow T_3$. Il y a un cycle, dont H n'est pas \textit{s\'erialisable}.
\end{solution}


\item Donner l'histoire finale par application de l'algorithme de verrouillage \`a deux phases. 
Donner le d\'etail du d\'eroulement de l'algorithme.
 \texttt{(2 points)}

\begin{solution}
$T_1$  est mis en attente au moment où $w_1[x]$ est soumis ($x$ a un  verrou partagé avec $T_2$)
$T_2$  est mis en attente au moment où $w_2[y]$ est soumis ($x$ a un  verrou partagé avec $T_3$). 
$T_3$ finit alors de s'exécuter, $T_2$ est libérée et termine, puis $T_1$.
\end{solution}


\end{enumerate}
